\section{Key Performance Indicators}

%___________________________________________________________________
\begin{frame}{KPIs: Stabilität}
    \framesubtitle{Struktur \& Stabilität im Institutsvergleich}

    \begin{columns}[T]
        \begin{column}{0.48\textwidth}
            \onslide<2->{
                \kpi{Wirtsch. Eigenkapitalquote}{\qty{78,1}{\percent}}%
                    {(Eigenkapital + Sonderposten) / Bilanzsumme}%
                    {Exzellenter Wert, Sonderposten wirken wie risikofreies Eigenkapital. Nominale EQ ist mit \qty{14,0}{\percent} deutlich niedriger.}
            }
                
            \onslide<4->{
                \kpi{Anlagenintensität}{\qty{61,5}{\percent}}%
                    {Anlagevermögen / Bilanzsumme}%
                    {Strukturelle Besonderheit: Extrem hoch im Vergleich zu IT-Instituten, aber absolut typisch für die hardwarenahe Halbleiterforschung.}
            }
        \end{column}

        \begin{column}{0.48\textwidth}
            \onslide<3->{
                \kpi{Working Capital}{\qty{20,9}{\mega\euro}}%
                    {Umlaufvermögen -- Verbindlichkeiten}%
                    {Sehr komfortabler Puffer. Sichert die vorübergehende Vorfinanzierung laufender Projekte.}
            }

            \onslide<5->{
                \kpi{Investitionsquote}{\qty{27,4}{\percent}}%
                    {Zuweisung Zuschüsse / Anlagevermögen}%
                    {Weit über dem Industriedurchschnitt durch Modernisierung der Reinraumanlagen.}
            }
        \end{column}
    \end{columns}

\end{frame}


%___________________________________________________________________
\begin{frame}{KPIs: Leistung}
    \framesubtitle{Operative Exzellenz und Drittmittelstärke}

    \begin{columns}[T]
        \begin{column}{0.48\textwidth}
            \onslide<2->{
                \kpi{Drittmittelquote}{\qty{29,7}{\percent}}%
                    {Erlöse / Verfügbare Erträge}%
                    {Genau im Mittelwert der Leibniz-Gemeinschaft. Belegt die starke Brückenfunktion zwischen Grundlagenforschung und Industrie.}
            }

            \onslide<4->{
                \kpi{Arbeitsproduktivität}{\qty{85,4}{\kilo\euro} / FTE}%
                    {Erlöse / wissenschaftliche FTE}%
                    {Hohe pro-Kopf-Einwerbung bei 185 wissenschaftlichen Vollzeitäquivalenten.}
            }
        \end{column}

        \begin{column}{0.48\textwidth}
            \onslide<3->{
                \kpi{Operativer Cash-Flow}{\qty{26,6}{\mega\euro}}%
                    {Jahresüberschuss + Abschreibungen + Zuweisung zu Sonderposten}%
                    {Enormer Finanzierungsstrom, der jedoch komplett zweckgebunden ist und direkt in die Reinvestition fließt.}
            }

            \onslide<5->{
                \kpi{Personalaufwandsquote}{\qty{49,4}{\percent}}%
                    {Personalaufwand / Verfügbare Erträge}%
                    {Hohe Sachkosten aufgrund der energieintensiven Reinraum-Infrastruktur.}
            }
        \end{column}
    \end{columns}

\end{frame}